% === Begin Label Data ===
% ID: 881
% 难度星级: 
% 标签: 
% 备注: 新定义
% 组卷引用次数: 0
% === End  Label Data ===

\begin{problem}{2019}{G}{北京卷}{20}{数列}
已知数列 $\{a_n\}$, 从中选取第 $i_1$ 项、第 $i_2$ 项、$\dots$、第 $i_m$ 项 ($i_1 < i_2 < \dots < i_m$), 若 $a_{i_1} < a_{i_2} < \dots < a_{i_m}$, 则称新数列 $a_{i_1}, a_{i_2}, \dots, a_{i_m}$ 为 $\{a_n\}$ 的长度为 $m$ 的递增子列. 规定: 数列 $\{a_n\}$ 的任意一项都是 $\{a_n\}$ 的长度为 $1$ 的递增子列.

（I）写出数列 $1,8,3,7,5,6,9$ 的一个长度为 $4$ 的递增子列;

（II）已知数列 $\{a_n\}$ 的长度为 $p$ 的递增子列的末项的最小值为 $a_{m_0}$, 长度为 $q$ 的递增子列的末项的最小值为 $a_{n_0}$. 若 $p < q$, 求证: $a_{m_0} < a_{n_0}$;

（III）设无穷数列 $\{a_n\}$ 的各项均为正整数, 且任意两项均不相等. 若 $\{a_n\}$ 的长度为 $s$ 的递增子列末项的最小值为 $2s-1$, 且长度为 $s$ 末项为 $2s-1$ 的递增子列恰有 $2^{s-1}$ 个 ($s=1,2,\dots$), 求数列 $\{a_n\}$ 的通项公式.
\end{problem}

\begin{answer}
（I）$1,3,5,6$ (答案不唯一)

（III）$a_n = \begin{cases} 2k-1, & n=2k \\ 2k, & n=2k-1 \end{cases}$
\end{answer}

\begin{solutions}
（II）设长度为 $q$ 末项为 $a_{n_0}$ 的一个递增子列为 $a_{r_1}, a_{r_2}, \dots, a_{r_{q-1}}, a_{n_0}$.
由 $p < q$, 得 $a_{r_p} \leqslant a_{r_{q-1}} < a_{n_0}$.
因为 $\{a_n\}$ 的长度为 $p$ 的递增子列末项的最小值为 $a_{m_0}$,
又 $a_{r_1}, a_{r_2}, \dots, a_{r_p}$ 是 $\{a_n\}$ 的长度为 $p$ 的递增子列,
所以 $a_{m_0} \leqslant a_{r_p}$, 所以 $a_{m_0} < a_{n_0}$.

（III）简易证法, 不涉及数学归纳法, 逻辑通顺可得满分:
由题设易得, $1^*$所有正奇数都是 $\{a_n\}$ 中的项.
先证明 $2^*$若 $2m$ 是 $\{a_n\}$ 中的项, 则 $2m$ 必排在 $2m-1$ 之前 ($m$ 为正整数).
假设 $2m$ 排在 $2m-1$ 之后.
设 $a_{p_1}, a_{p_2}, \dots, a_{p_{m-1}}, 2m-1$ 是数列 $\{a_n\}$ 的长度为 $m$ 末项为 $2m-1$ 的递增子列, 则
$a_{p_1}, a_{p_2}, \dots, a_{p_{m-1}}, 2m-1, 2m$ 是数列 $\{a_n\}$ 的长度为 $m+1$ 末项为 $2m$ 的递增子列, 矛盾;
再证明 $3^*$: 所有正偶数都是 $\{a_n\}$ 中的项.

假设存在正偶数不是 $\{a_n\}$ 中的项, 设不在 $\{a_n\}$ 中的最小的正偶数为 $2m$. 因为 $2k$ 排在 $2k-1$ 之前 ($k=1,2,\dots,m-1$), 所以 $2k$ 和 $2k-1$ 不可能在 $\{a_n\}$ 中的同一个递增子列中. 又 $\{a_n\}$ 中不超过 $2m+1$ 的数为 $1,2,\dots,2m-2,2m-1,2m+1$, 所以 $\{a_n\}$ 的长度为 $m+1$ 且末项为 $2m+1$ 的递增子列个数至多为 $\underbrace{2 \times 2 \times 2 \times \dots \times 2}_{(m-1)\text{个}} \times 1 \times 1 = 2^{m-1} < 2^m$, 矛盾;

最后证明 $4^*$: $2m$ 排在 $2m-3$ 之后 ($m \geqslant 2$ 为整数).

假设存在 $2m (m \geqslant 2)$, 使得 $2m$ 排在 $2m-3$ 之前,

则 $\{a_n\}$ 的长度为 $m+1$ 且末项为 $2m+1$ 的递增子列的个数小于 $2^m$, 矛盾;

综上, 数列 $\{a_n\}$ 只可能为 $2,1,4,3,\dots,2m-3,2m,2m-1,\dots$.

经验证, 数列 $2,1,4,3,\dots,2m-3,2m,2m-1,\dots$ 符合条件.

所以 $a_n = \begin{cases} 2k-1, & n=2k \\ 2k, & n=2k-1 \end{cases}$

（III）标准证法: 先进行简单的数学实验:
首先考虑 $s=1$ 时,
长度为 $1$ 的递增子列末项的最小值为 $1$, 且长度为 $1$ 末项为 $1$ 的递增子列恰有 $1$ 个,
显然可得 $1$ 在数列中;

考虑 $s=2$ 时,
长度为 $2$ 的递增子列末项的最小值为 $3$, 且长度为 $2$ 末项为 $3$ 的递增子列恰有 $2$ 个,
显然可得 $3$ 在数列中,
$2$ 在数列中 (若 $2$ 不在数列中只有 $13$ 一个递增子列, 不够)
(若 $2$ 在 $1$ 的后面, 则会有 $12$, 末项最小值过小)

考虑 $s=3$ 时,
长度为 $3$ 的递增子列末项的最小值为 $5$, 且长度为 $3$ 末项为 $5$ 的递增子列恰有 $4$ 个,
显然可得 $5$ 在数列中, 且在 $213$ 的后面,
$4$ 在数列中 (若 $4$ 不在数列中只有 $135, 235$ 两个递增子列, 不够)
$4$ 在 $3$ 的前面 (若 $4$ 在 $3$ 的后面, 则会有 $134$, 末项最小值过小)
$4$ 在 $1$ 后面 (若 $4$ 在 $1$ 的前面, 则只会有 $241, 235, 135$ 三个递增子列, 不够)

考虑 $s=4$ 时,
长度为 $4$ 的递增子列末项的最小值为 $7$, 且长度为 $4$ 末项为 $7$ 的递增子列恰有 $8$ 个,
可得 $7$ 在数列中, 在 $21435$ 的后面, $6$ 在数列中且在 $5$ 的前面, $3$ 的后面, 原因同上;
所以考虑 $a_n = \begin{cases} 2k-1, & n=2k \\ 2k, & n=2k-1 \end{cases}$

下用数学归纳法尝试证明:

$n=1,2,3$ 时已经完成证明;

假设 $n=2i-2, 2i-1$ 时成立, 则前出现顺序为 $2,1,4,3,6,5,\dots,2i-2,2i-3,2i-1$

下证明 $n=2i, 2i+1$ 时成立:
考虑 $2i+1$, 显然出现在数列中, 考虑位置:
末项为 $2i+1$ 的长为 $i+1$ 的递增数列, 其前 $i$ 项必定分别取自集合 $\{2,1\}, \{4,3\}, \dots, \{2i, 2i-1\}$,
共 $2^i$ 个, 所以 $2k+1$ 必然在 $2k-1$ 之后;
考虑 $2i$, 末项为 $2i+1$ 的长为 $i+1$ 的递增数列需有 $2^i$ 个
其前 $i$ 项必定分别取自集合 $\{2,1\}, \{4,3\}, \dots, \{2i-2, 2i-3\}, 2i-1$,
若它未出现, 共 $2^{i-1}$ 个, 矛盾;

若它出现在 $2i-1$ 之后, 则可构造末项为 $2i$ 的长为 $i+1$ 的递增数列, 矛盾;
若它出现在 $2i-3$ 之前, 则末项为 $2i+1$ 的长为 $i+1$ 的递增数列,
其前 $i$ 项必定分别取自集合 $\{2,1\}, \{4,3\}, \dots, \{2i, 2i-1\}$, 至多 $2^i$ 个,
但 $1,3,5,\dots,2i-1, 2i, 2i+1$, 这个递增数列取不到, 故没有 $2^{i+1}$ 个所求数列, 矛盾;
则前 $2i+1$ 项出现顺序为只能 $2,1,4,3,6,5,\dots,2i-2,2i-3,2i,2i-1,2i+1$
故 $n=2i, 2i+1$ 时成立, $a_n = \begin{cases} 2k-1, & n=2k \\ 2k, & n=2k-1 \end{cases}$. 证明完成.
标黄处的数学实验与数学归纳法的作用可仔细体会.
\end{solutions}